\documentclass{report}
\usepackage{amsmath,amssymb}
\setlength{\parindent}{0mm}
\setlength{\parskip}{1em}
\begin{document}
\begin{center}
\rule{6in}{1pt} \
{\large Christopher Leibs \\
{\bf Least-Squares Finite Element Method and Nested Iteration for Spatially Variable Temperatures in a Two-Fluid Plasma}}

526 UCB \\ University of Colorado \\ Boulder CO \\ 80309
\\
{\tt leibs@colorado.edu}\\
Thomas Manteuffel\end{center}

Previously, the authors introduced a scalable solution technique for an
electromagnetic, two-fluid plasma (TFP) model that utilized least-squares
finite elements, nested iteration, and adaptive mesh refinement [1]. The
system under investigation used a naive adiabatic (constant) temperature
model. Such an assumption was key in designing a scaling and modification
of the TFP model to produce optimal finite element performance, and
linear systems amenable to solution with Algebraic Multigrid (AMG). To
efficiently focus computational resources, an adaptive mesh refinement
scheme was used.

In this talk, we discuss the recent exploration of spatially variable
temperature models. The motivation is twofold: (1) many fluid-based
temperature closures can simply be lagged in nonlinear iteration or
nested iteration with little overhead, and (2) the ability to provide a
kinetic-based closure becomes viable. Without loss of generality, we
focus on the former and consider how the previous solution technique is
modified in order to handle such variability while maintaining
scalability and efficiency.

[1] C. A. Leibs, T. A. Manteuffel, SIAM Journal on Scientific Computing,
37(5):S314-S333, 2015.


\end{document}
